% 第九卷第3章：全注意力的键值缓存与递归汇聚的状态更新对照（跨栏，约150 mm宽）。
\begin{tikzpicture}[
  x=1mm,y=1mm,
  figure picture,
  cell/.style={rectangle,draw=FigureInput,fill=FigureInputFill,line width=.6pt,
    minimum width=10mm,minimum height=6mm,inner sep=0pt,font=\figurefont\footnotesize},
  vcell/.style={cell,draw=FigureModel,fill=FigureModelFill},
  qnode/.style={figure node,draw=FigureOutput,fill=FigureOutputFill,minimum height=7mm,
    minimum width=10mm},
  state/.style={figure node,draw=FigureModel,fill=FigureModelFill,minimum width=16mm,
    minimum height=14mm},
  op/.style={figure node,minimum height=9mm,minimum width=26mm},
  dop/.style={op,dash pattern=on 2pt off 1.5pt},
  plus/.style={circle,draw=FigureStroke,line width=.75pt,minimum size=5mm,inner sep=0pt,
    font=\figurefont\footnotesize},
  lab/.style={figure note,anchor=west}
]
  % ---------- 上方：全注意力的键值缓存 ----------
  \node[figure heading,anchor=west] at (0,30) {全注意力：逐项追加键值缓存};
  \foreach \i/\name in {0/1,1/2,2/3} {
    \node[cell] (k\name) at ({8+\i*10},20) {$\bm k_{\name}$};
    \node[vcell] (v\name) at ({8+\i*10},13) {$\bm v_{\name}$};
  }
  \node[font=\figurefont\footnotesize] at (38,16.5) {$\cdots$};
  \node[cell] (kt1) at (48,20) {$\bm k_{t-1}$};
  \node[vcell] (vt1) at (48,13) {$\bm v_{t-1}$};
  \node[cell,dash pattern=on 2pt off 1.5pt] (kt) at (58,20) {$\bm k_{t}$};
  \node[vcell,dash pattern=on 2pt off 1.5pt] (vt) at (58,13) {$\bm v_{t}$};
  \draw[figure link] (3,8) -- (63,8);
  \draw[figure link] (3,7) -- (3,9);
  \draw[figure link] (63,7) -- (63,9);
  \node[figure note] at (33,4.5) {缓存长度$t$，随生成线性增长};
  \node[qnode] (q) at (82,16.5) {$\bm q_t$};
  \foreach \n in {k1,k2,k3,kt1,kt} {
    \draw[figure link,densely dotted] (\n.north) -- (\n.north |- 0,26);
  }
  \draw[figure link,densely dotted] (k1.north |- 0,26) -- (70,26) |- (q.west);
  \node[lab] at (90,19) {与全部$t$个键配对打分，};
  \node[lab] at (90,14) {再对$t$个值加权求和};

  % 分隔线
  \draw[draw=FigureGrid,line width=.5pt] (0,0) -- (150,0);

  % ---------- 下方：递归汇聚的固定大小状态 ----------
  \node[figure heading,anchor=west] at (0,-6) {递归汇聚：更新固定大小的状态矩阵};
  \node[state] (sold) at (10,-24) {$\bm S_{t-1}$\\[-1pt]{\scriptsize $d_v\times d_k$}};
  \node[dop] (decay) at (45,-24) {衰减或擦除\\{\scriptsize 修改旧状态}};
  \node[plus] (add) at (72,-24) {$+$};
  \node[op,draw=FigureInput,fill=FigureInputFill] (write) at (72,-42) {写入外积\\$\beta_t\bm v_t\bm k_t^{\top}$};
  \node[state] (snew) at (98,-24) {$\bm S_{t}$\\[-1pt]{\scriptsize $d_v\times d_k$}};
  \node[qnode] (q2) at (98,-45) {$\bm q_t$};
  \node[figure node,draw=FigureOutput,fill=FigureOutputFill,minimum height=9mm,minimum width=22mm] (out) at (132,-24) {读出\\$\bm o_t=\bm S_t\bm q_t$};

  \draw[figure flow] (sold.east) -- (decay.west);
  \draw[figure flow] (decay.east) -- (add.west);
  \draw[figure flow] (write.north) -- (add.south);
  \draw[figure flow] (add.east) -- (snew.west);
  \draw[figure flow] (snew.east) -- (out.west);
  \draw[figure flow] (q2.east) -| (out.south);
  \draw[figure flow,draw=FigureMuted,dash pattern=on 2pt off 1.5pt]
    (snew.north) -- ++(0,6) -| (sold.north);
  \node[figure note] at (54,-13.5) {下一步作为旧状态};
\end{tikzpicture}
